% =====================================================================
\chapter*{Sources et références}
\addcontentsline{toc}{chapter}{Sources et références}
\markboth{Sources et références}{Sources et références}
% =====================================================================

\section*{Sources de données utilisées}

\begin{itemize}[leftmargin=1.4em]
  \item \textbf{Géorisques} — portail national des risques naturels et
        technologiques, Ministère de la Transition écologique / BRGM.
        Données distribuées sous Licence Ouverte 2.0 (Etalab).
        \url{https://www.georisques.gouv.fr}
  \item \textbf{BDNB} — Base de Données Nationale des Bâtiments, Centre
        Scientifique et Technique du Bâtiment (CSTB). Licence Ouverte 2.0.
        \url{https://api.bdnb.io}
  \item \textbf{Géoplateforme} — service de géocodage, Institut Géographique
        National (IGN). \url{https://data.geopf.fr/geocodage}
  \item \textbf{Sandre} — portail national d'accès aux données sur l'eau.
        \url{https://www.sandre.eaufrance.fr}
  \item \textbf{Open-Meteo} — API météorologique ouverte, incluant le modèle
        hydrologique GloFAS. \url{https://open-meteo.com}
  \item \textbf{Panoramax} — service public de photographie de terrain
        contributive. \url{https://panoramax.fr}
  \item \textbf{Copernicus Climate Data Store} — projections climatiques
        européennes, ECMWF. Jeu de données \emph{Climate indicators for Europe
        from 1940 to 2100}. \url{https://cds.climate.copernicus.eu}
  \item \textbf{Vigicrues} — service public de vigilance crues.
        \url{https://www.vigicrues.gouv.fr}
  \item \textbf{API Carto / BAN} — Base Adresse Nationale.
        \url{https://adresse.data.gouv.fr}
\end{itemize}

\section*{Cadres réglementaires}

\begin{itemize}[leftmargin=1.4em]
  \item \textbf{Licence Ouverte 2.0} (Etalab) — cadre juridique de
        réutilisation des données publiques françaises. Réutilisation
        commerciale autorisée, obligation d'attribution.
        \url{https://www.etalab.gouv.fr/licence-ouverte-open-licence}
  \item \textbf{Directive Inondation} — cadre européen ayant conduit à la
        définition des Territoires à Risque Important d'inondation (TRI).
  \item \textbf{Régime CatNat} — dispositif français d'indemnisation des
        catastrophes naturelles, article L125-1 du Code des assurances.
  \item \textbf{Solvabilité II} — directive européenne 2009/138/CE encadrant la
        solvabilité des organismes d'assurance et de réassurance.
\end{itemize}

\section*{Standards techniques}

\begin{itemize}[leftmargin=1.4em]
  \item \textbf{OGC Web Feature Service 2.0} — standard d'échange de données
        géographiques vectorielles.
  \item \textbf{OGC Geography Markup Language 3.2} — format d'encodage des
        géométries.
  \item \textbf{Mapbox GL JS} — bibliothèque de rendu cartographique.
  \item \textbf{FastAPI} — cadre de développement d'API Python.
  \item \textbf{Pydantic v2} — validation de données Python.
\end{itemize}

\section*{Documents internes au projet}

\begin{itemize}[leftmargin=1.4em]
  \item \fichier{constitution.md} — règles non négociables, version 0.2,
        amendée le 14 septembre 2026.
  \item \fichier{docs/workflow/01-brainstorm-notes.md} — journal des décisions
        de la phase 1.
  \item \fichier{docs/workflow/concept.md} — concept et questions ouvertes
        résolues (OQ-1 à OQ-6).
  \item \fichier{docs/workflow/spec.md} — spécification fonctionnelle,
        28 exigences \frref{01} à \frref{28}.
  \item \fichier{docs/workflow/plan.md} — plan d'exécution, tâches T001 à T018.
  \item \fichier{docs/workflow/scenario-plan.md} — plan du moteur de scénarios,
        tâches \scn{001} à \scn{050}.
  \item \fichier{docs/workflow/vfx-spec.md} et \fichier{vfx-plan.md} —
        spécification et plan du mode de rendu cinématique.
  \item \fichier{AGENTS.md} — conventions, invariants acquis et pièges
        techniques.
  \item \fichier{handoff.md} — décisions architecturales en attente.
  \item \fichier{roadmapproduction.md} — feuille de route vers la production.
  \item \fichier{README.md} — documentation d'installation, d'usage et de
        déploiement.
\end{itemize}
